%% =====================================================================
%%  CAPITULO 9 - DISCUSSAO
%% =====================================================================

\chapter{Discussão}
\label{cap:discussao}

Este capítulo interpreta os resultados, confronta-os com o problema e a
hipótese, identifica as lacunas que a própria verificação revelou e retoma as
ameaças à validade. A seção final discute em que medida os achados são
generalizáveis.

\section{Retomada do problema e da hipótese}
\label{sec:disc-problema}

O problema enunciado na Seção~\ref{sec:problema} tinha três componentes:
rastreabilidade a requisito, verificabilidade por critério e sustentação por
evidência reproduzível. Os três foram examinados, com resultados desiguais.

A \textbf{verificabilidade por critério} foi o componente mais bem-sucedido. Os
doze critérios declarados foram satisfeitos, e cada um foi verificado por
observação externa, sem consultar o autor e sem instrumentar o produto. A
consequência relevante é metodológica: como os critérios foram redigidos de forma
que um observador externo possa decidir o resultado, verificar deixou de exigir
acesso ao raciocínio de quem implementou.

A \textbf{sustentação por evidência reproduzível} também se confirmou. A
inspeção estática é reproduzível por comando; a execução observada é reproduzível
por automação; a medição de intervalo é reproduzível por agente de teste que
observa a superfície de desenho. Nenhuma conclusão deste trabalho depende de
lembrança ou de registro narrativo.

A \textbf{rastreabilidade a requisito} foi o componente que falhou parcialmente.
A cobertura de 54,2\,\% significa que quase metade dos requisitos não possui
critério de aceitação associado. A cadeia requisito--critério--evidência existe,
mas apenas para parte do conjunto especificado. A hipótese, portanto, é
\emph{parcialmente} sustentada: o processo produz entregas verificáveis, mas não
garante que a verificabilidade cubra a especificação inteira.

\section{A assimetria entre requisitos funcionais e não funcionais}
\label{sec:disc-assimetria}

O achado mais relevante do trabalho é a diferença de cobertura entre as classes
de requisito: 73,3\,\% para funcionais e 22,2\,\% para não funcionais. A
magnitude da diferença descarta a explicação por acaso.

A explicação é estrutural. Um requisito funcional descreve um comportamento
observável, e comportamento observável tem forma de teste quase natural: dado um
estado, ao executar uma ação, espera-se uma consequência. Requisito não funcional
enunciado como atributo --- ``layout moderno e limpo'', ``tema escuro de estilo
retro-arcade'' --- não descreve consequência observável; descreve uma
característica percebida. Sem um passo intermediário de conversão em métrica, ele
não gera critério, e sem critério não entra em portão algum.

Esse passo intermediário não é trivial e não foi executado. A
Tabela~\ref{tab:mapeamento-iso} mostra o sintoma: três dos nove requisitos não
funcionais recaem em características de qualidade para as quais nenhuma métrica
foi declarada. O requisito existe, a característica de qualidade foi
identificada, mas o elo que os liga a um teste não foi construído.

A conclusão prática é que a exigência de critério de aceitação deve ser
\emph{universal}, não apenas aplicada aos requisitos que naturalmente se
convertem em teste. Um requisito sem critério é um requisito não verificável, e um
requisito não verificável é indistinguível de uma intenção. O framework precisa
de um portão que rejeite requisito órfão de critério, ou que force sua
conversão em métrica ou seu rebaixamento explícito a diretriz não verificável.

\subsection{A lacuna do portão de intenção}
\label{sec:disc-lacunas-do-gate}

O portão G1 exige que existam um requisito identificado e um critério de
aceitação antes do avanço para a implementação. A redação do portão, entretanto,
é satisfeita pela existência de \emph{ao menos} um par requisito--critério. Ela
não impõe que todo requisito declarado possua critério correspondente.

A consequência é assimétrica e foi medida: o portão deixa passar uma
especificação em que onze requisitos funcionais e dois não funcionais têm
critério, e dez requisitos não têm nenhum. A verificação subsequente de doze
critérios produz a impressão de cobertura completa, porque todos os critérios
declarados foram satisfeitos. A lacuna permanece invisível até que alguém
confronte explicitamente a contagem de requisitos com a contagem de critérios ---
procedimento que, no caso deste trabalho, só foi executado na etapa de
consolidação da matriz de rastreabilidade.

Duas correções são possíveis, e elas não são equivalentes. A primeira é tornar o
portão mais estrito: bloquear o avanço enquanto houver requisito sem critério. A
segunda é tornar o portão mais informativo: exigir que todo requisito sem
critério seja explicitamente rebaixado a ``diretriz não verificável neste
ciclo''. A segunda é preferível, porque existem requisitos legitimamente não
verificáveis por teste automatizado --- aparência visual, por exemplo --- e
bloquear a implementação por causa deles tornaria o processo inaplicável. O
rebaixamento explícito resolve o problema real, que não é a ausência de teste, e
sim a ausência de \emph{consciência} da ausência de teste.

Formalizando: seja $R$ o conjunto de requisitos e $R_c \subseteq R$ o subconjunto
com critério de aceitação dedicado. O portão atual exige $|R_c| \geq 1$. A
correção proposta exige que $R \setminus R_c$ esteja declarado como conjunto de
diretrizes não verificáveis neste ciclo, e não que $R \setminus R_c = \varnothing$.
Neste estudo de caso, $|R| = 24$ e $|R_c| = 13$, de modo que
$|R \setminus R_c| = 11$: dez requisitos sem critério algum e um coberto apenas
parcialmente. Esse é o conjunto que hoje não está declarado como não verificável.

\section{Valor dos casos de borda e das decisões arquiteturais}
\label{sec:disc-bordas}

Dois resultados empíricos sustentam a tese de que declarar casos de borda e
registrar decisões não é cerimônia documental.

O primeiro é o estado de vitória. O caso de borda ``grade completamente
ocupada'' foi declarado na especificação antes da implementação. Ao implementar
o sorteio do alvo, tornou-se evidente que essa situação exigia um estado de jogo
que não existia na modelagem inicial. O caso de borda \emph{criou} um estado
(Figura~\ref{fig:fsm}). Sem ele, o sorteio retornaria valor inválido em uma
situação alcançável, e o defeito só apareceria quando alguém completasse o
tabuleiro --- provavelmente nunca, em uso normal.

O segundo é o comportamento da cauda. O caso de borda~5 e a decisão ADR-004
combinam-se para produzir um comportamento observável e não intuitivo: girando
em quadrado com comprimento quatro o jogador sobrevive; com comprimento cinco,
colide consigo mesmo (Seção~\ref{sec:res-extras}, item~4). Dois implementadores
igualmente competentes, partindo da mesma especificação de requisitos e sem o
registro de ADR-004, produziriam comportamentos diferentes. A decisão
arquitetural é o que torna o resultado determinístico. Isso responde ao argumento
de que registrar decisões é documentação supérflua: nesse caso, o registro
elimina uma ambiguidade real de implementação.

\section{Custo do processo}
\label{sec:disc-custo}

O processo tem custo, e ele deve ser declarado. A disciplina exigida produziu
cinco documentos de especificação e projeto antes da primeira linha de código:
especificação, decisões arquiteturais, tarefas, evidências e controle de estado.
Para um produto de 1\,063 linhas, essa proporção é alta.

Três observações qualificam o custo.

\textbf{O custo é antecipado, não adicionado.} A declaração prévia de casos de
borda evitou defeitos que seriam descobertos tardiamente, quando o custo de
correção é maior (Equação~(\ref{eq:custo-fase})). O estado de vitória é o exemplo
concreto: descobri-lo em produção custaria mais do que especificá-lo.

\textbf{O custo é amortizado na verificação.} Critérios redigidos antes da
implementação produziram verificação executável por terceiros. Critérios
redigidos depois teriam exigido reinterpretação do autor a cada verificação.

\textbf{O custo é proporcional ao risco, não ao tamanho.} O fluxo adaptativo
(Seção~\ref{sec:fw-fluxo}) existe precisamente para evitar aplicar a cadeia
completa a alterações pequenas. Neste estudo de caso, a trilha aplicada foi a de
recurso médio com elevação a amplo, decisão coerente com o fato de o alvo ter
sido especificado do zero.

Não foi possível quantificar o custo em tempo ou em tokens, porque a execução não
foi instrumentada (P-12, P-13). A avaliação acima é qualitativa e não substitui a
medição ausente.

\section{Ameaças à validade revisitadas}
\label{sec:disc-ameacas}

As sete ameaças do Quadro~\ref{qua:ameacas} são retomadas aqui quanto ao efeito
sobre as conclusões.

\textbf{AV-1, ausência de grupo de controle, é a ameaça mais séria.} Sem um grupo
que executasse o mesmo alvo sem o processo, não é possível atribuir os doze
critérios satisfeitos ao processo. Uma explicação alternativa é que o alvo é
simples e que um implementador competente obteria o mesmo resultado sem
especificação formal. O trabalho não refuta essa alternativa. O que ele
demonstra é mais modesto e ainda assim relevante: que o processo torna o
resultado \emph{verificável por terceiros sem acesso ao histórico}, propriedade
que não foi medida no grupo alternativo justamente porque esse grupo não existiu.

\textbf{AV-6, amostra única, limita a medição de velocidade.} O intervalo medido
é coerente com o previsto, mas coerência em uma execução não é confirmação
estatística. Um resultado mais forte exigiria repetição e intervalo de confiança.

\textbf{AV-3, autoria coincidente, é mitigada pela natureza da evidência.} A
verificação foi produzida por execução observada e inspeção textual, não por
julgamento do autor. Um terceiro pode repetir os comandos e obter os mesmos
resultados. Isso não elimina o risco de confirmação na escolha dos critérios, mas
remove o risco de confirmação na produção da evidência.

\textbf{AV-4, a métrica de rastreabilidade mede existência e não qualidade.} Um
critério mal redigido contaria como cobertura. A ameaça é real e reduz o valor do
número de 54,2\,\%. Ela não afeta, entretanto, o achado principal, que é a
\emph{assimetria} entre as classes: se critérios ruins inflassem artificialmente
a cobertura, a cobertura de requisitos não funcionais seria ainda menor, e a
assimetria, mais acentuada.

\textbf{AV-2 e AV-5} permanecem nos limites declarados: generalização restrita a
domínios de complexidade comparável e viés positivo nos valores absolutos de
intervalo.

\section{Generalização}
\label{sec:disc-generalizacao}

A generalização de um estudo de caso único é analítica, não estatística
(\citeonline{yin2018}): o resultado não se estende a uma população, mas
confirma ou refuta proposições teóricas. Três proposições recebem suporte
limitado dos dados.

\textbf{Proposição 1.} Requisitos funcionais convertem-se em critério de
aceitação com mais facilidade que requisitos não funcionais. Sustentada por
evidência direta neste caso (73,3\,\% contra 22,2\,\%). A proposição é plausível
em qualquer domínio de software; a magnitude específica não é transferível.

\textbf{Proposição 2.} Declarar casos de borda antes da implementação altera o
projeto, não apenas os testes. Sustentada por evidência direta (o estado de
vitória). A proposição decorre diretamente dos dados e não depende do domínio.

\textbf{Proposição 3.} Evidência produzida por observação externa ao produto é
reproduzível por terceiros sem instrumentar o código. Sustentada pela construção
metodológica deste trabalho, não por comparação.

\textbf{O que não é transferível.} O domínio do estudo de caso tem estado
discreto, regras determinísticas e ausência de concorrência e de integração
externa. Em domínios com estado contínuo, tempo real estrito, sistema distribuído
ou dependência de hardware, a conversão de requisitos em critério é
substancialmente mais difícil, e o processo exigiria níveis de evidência que este
trabalho não alcança. O próprio material de origem do framework previa quatro
níveis de evidência --- teste em máquina hospedeira, simulação, bancada e
integração com hardware --- e este estudo permaneceu no nível mais baixo. A
aplicação a domínios de nível mais alto permanece como questão aberta.

\section{Comparação com a literatura}
\label{sec:disc-literatura}

Três convergências merecem registro.

A ênfase de \citeonline{boehm1984} na verificação precoce de requisitos e projeto
encontra apoio no achado da Seção~\ref{sec:disc-bordas}: a declaração prévia de
casos de borda alterou a modelagem de estados antes que houvesse código.

A tese de \citeonline{adzic2011} de que o exemplo concreto serve
simultaneamente como requisito, teste e documentação é confirmada pelo formato
dos critérios usados neste trabalho, redigidos em forma de situação, ação e
consequência esperada.

A observação de \citeonline{myers2011} de que um caso de teste que não pode
falhar não tem valor é o critério pelo qual se identifica a lacuna dos requisitos
não funcionais: os sete requisitos sem critério não possuem caso de teste capaz de
reprová-los, e por isso não entraram em nenhum portão.

Em relação a \citeonline{hong2024metagpt}, \citeonline{qian2024chatdev} e
\citeonline{wu2023autogen}, a diferença de ênfase apontada no
Capítulo~\ref{cap:relacionados} é confirmada por este estudo: aqueles trabalhos
maximizam a qualidade do artefato gerado pela coordenação de papéis; este trabalho
deslocou o foco para os artefatos que tornam a geração auditável. Os dois enfoques
são complementares, não concorrentes.

\section{Implicações práticas}
\label{sec:disc-implicacoes}

Quatro implicações decorrem dos resultados.

\textbf{Para o processo.} A exigência de critério deve ser universal e verificada
por portão. Requisito sem critério deve ser bloqueado ou explicitamente
rebaixado a diretriz não verificável.

\textbf{Para a especificação.} Requisito não funcional deve ser redigido já com
sua métrica, ou com a declaração explícita de que não é verificável neste ciclo.
A declaração explícita é preferível à omissão, porque torna a lacuna visível.

\textbf{Para a verificação.} Observar o produto por fora, sem instrumentá-lo,
resolve o conflito entre testar e alterar o objeto de teste. Em interfaces
gráficas discretas, a amostragem da superfície de desenho é uma técnica viável e
de baixo custo.

\textbf{Para a condução por agentes.} As condições de parada cumpriram função
real: a divergência entre a documentação do repositório e seu conteúdo efetivo
foi detectada e registrada em vez de resolvida por suposição silenciosa. Sem uma
condição de parada explícita, a resposta mais provável teria sido produzir
documentação sobre um sistema inexistente.
